与插入时或阴道口的\textbf{浅层性交痛}不同，\textbf{深层性交痛}指的是在\textbf{插入较深、顶到盆腔深处时}出现的疼痛，多与盆腔内的病变有关，需要妇科检查才能明确原因。

\vspace{0.5em}
\noindent\textbf{1. 常见病因}

\begin{itemize}
	\item \textbf{子宫内膜异位症}：内膜组织长在子宫外，常累及子宫骶韧带、直肠阴道隔与卵巢。典型表现为\textbf{深部性交痛 + 进行性痛经 + 慢性盆腔痛}的"三联征"，疼痛常在月经前后加重。
	\item \textbf{盆腔炎性疾病（PID）及其后遗粘连}：慢性盆腔炎、输卵管卵巢粘连会固定盆腔器官，性交顶撞时引起牵扯痛；常伴下腹坠痛、白带异常。
	\item \textbf{卵巢囊肿 / 卵巢病变}：较大的囊肿在性交时受挤压或牵扯可致痛，若突发剧烈腹痛须警惕囊肿破裂或扭转（急症）。
	\item \textbf{子宫腺肌症、子宫后位固定}：子宫增大、活动度差时，深部性交痛常见。
	\item \textbf{术后粘连、盆底高张力}：盆腔手术、子宫内膜异位术后瘢痕，或盆底肌长期紧张，也会造成深部痛。
\end{itemize}

\noindent\textbf{2. 如何就医}

\begin{itemize}
	\item 尽量\textbf{描述清楚}：是插入时痛、顶到深处才痛，还是抽动中痛；是钝痛、刺痛还是牵扯痛；与月经周期、体位、排便的关系——这些线索能帮助医生判断。
	\item 常用检查：\textbf{妇科双合诊/三合诊}、\textbf{经阴道超声}，必要时 MRI、腹腔镜（子宫内膜异位症诊断的"金标准"）。
	\item \textbf{出现以下情况应尽快就医}：突发剧烈下腹痛、伴发热或异常出血、疼痛持续加重，或伴不孕。
\end{itemize}

\noindent\textbf{3. 处理原则}

\begin{itemize}
	\item \textbf{治本}：针对病因治疗——子宫内膜异位症可用药物（激素治疗）或手术；PID 需规范抗感染；卵巢囊肿按指征处理。
	\item \textbf{缓解当下}：性交时采用\textbf{女方掌控深度与角度}的体位（如女上位、侧卧位），放慢节奏、避免过深顶撞；充分唤起与润滑能显著减轻不适。
	\item \textbf{盆底康复}：对盆底高张力相关者，盆底物理治疗（放松训练、生物反馈）有帮助。
	\item \textbf{别硬扛}：深部性交痛往往\textbf{提示盆腔存在需要处理的问题}，把它当作身体发出的信号，而不是"忍忍就过去"的小事。
\end{itemize}

\begin{tcolorbox}[colback=pink!5!white,colframe=red!50!black,title=贴心小叮咛]
浅层痛多与"入口"有关（阴道痉挛、前庭痛），深层痛则多与\textbf{盆腔内部病变}有关。如果疼痛只在插得较深时出现、或伴进行性痛经与慢性盆腔痛，请尽早到妇科就诊——\textbf{明确病因，疼痛才有被真正解决的可能}。
\end{tcolorbox}
